% Nouveau EGA V (cote 25): archive cote 25, site page 181 (PDF page 182).
% typescript. Draft compared: none.
% AI-transcribed by Claude Opus 5.5 (claude-opus-5-5, adaptive thinking, effort high), Message Batches API, batch ega25b (msgbatch_01A7n5YuXubUPNK1zsepi8RG), 2026-11-13; not yet checked by eye.
\scanpage{25/181}
\mspage{\hand{8}}

d'une somme directe de formes $Q_1 \tikz[baseline=(sum.base)]{\node[draw,circle,handwritten,text=black,inner sep=0pt] (sum) {$\unsure{+}$};} Q_2$ , $Q_1$ de rang 1 \unsure{e|a}t $Q_2$ de rang pair.
En vertu de 16.8. on voit que $Q_1,Q_2$ sont de discriminant corrigé
$\hand{\neq} 0$, en vertu de ce qui précède il s'ensuit que $Q_2$ est isomorphe à la
forme standart, et d'autre part $Q_1$ est non nulle, donc isomorphe à la
forme standart de degré 1, donc Q est isomorphe à la forme standart, ce
qui prouve 16.11.1.

Il reste à prouver l'équivalence de (i) et (ii) dans 16.9. On peut
évidemment supposer que S est le spectre d'un corps alg clos, \var{\inserted[typed]{et on est}}{et on est}{insertion}{ega25_25_181-v0001} \var{\xout{d'autre part}}{}{typewriter strike-out}{ega25_25_181-v0002}
ramené à prouver dans ce cas \tikz[baseline=(lemma.base)]{\node[draw,ellipse,handwritten,text=black,inner sep=2pt] (lemma) {\shortstack[l]{\var{\inserted{le}}{le}{insertion}{ega25_25_181-v0003} \var{\inserted[typed]{\uline{Lemme} 16.11.2.}}{\uline{Lemme} 16.11.2.}{insertion}{ega25_25_181-v0004}\\\var{\deleted{\unsure{que}}}{}{deletion}{ega25_25_181-v0005} Q lisse \hand{$\Longleftrightarrow$} Q standart.}};}
% NOTE: The pen oval encloses the lemma title and its formula; its upper edge crosses the preceding typed line without cancelling that line.
En vertu de 16.2. on peut supposer que k est de car 2 et E est de rang
impair. \var{\deleted{\var{\xout{On sait déjà que (ii) implique que}}{}{typewriter strike-out}{ega25_25_181-v0007}}}{}{deletion}{ega25_25_181-v0006}
% NOTE: Tiny dark specks above the initial On and a small n-like remnant below it have no secure added-word or symbol reading.
 \var{\inserted[typed]{Lorsque}}{Lorsque}{insertion}{ega25_25_181-v0008} Q est isomorphe à la forme
standart, \var{\xout{\unsure{on pourrait}}}{}{typewriter strike-out}{ega25_25_181-v0009} on vérifie aussitôt, grâce à 16.3. , qu'elle
est lisse (le noyau de $\overset{\hand{\vee}}{E}$ \hand{$\longrightarrow$}Ev est de dimension 1, et Q n'est pas identi-%
quement nulle sur ce noyau, donc n'y admet que le zéro trivial); \var{\xout{donc}}{}{typewriter strike-out}{ega25_25_181-v0010}
\var{\inserted[typed]{D'autre part}}{D'autre part}{insertion}{ega25_25_181-v0011} \var{\deleted{\var{\xout{(ii) \unsure{$\Rightarrow$} (i). Or}}{}{typewriter strike-out}{ega25_25_181-v0013}}}{}{deletion}{ega25_25_181-v0012} toujours par 16.3. , \var{\inserted[typed]{Q lisse}}{Q lisse}{insertion}{ega25_25_181-v0014} \var{\deleted{(ii)}}{}{deletion}{ega25_25_181-v0015} implique que \var{\xout{le noyau}}{}{typewriter strike-out}{ega25_25_181-v0016} la restric-%
tion de Q au noyau N de Q n'admet quenle zéro trivial, ce qui évidemment
\var{\inserted[typed]{signifie}}{signifie}{insertion}{ega25_25_181-v0017} \var{\xout{implique}}{}{typewriter strike-out}{ega25_25_181-v0018} que N est de dim 1 et que $Q|N \hand{\neq} 0$. Or prenant un supplémentaire
de N, on voit aussitôt que cela implique que Q est isomorphe à la forme
standar\unsure{d|t}\hand{/}. Cela \var{\inserted[typed]{prouve 16.11.2. et}}{prouve 16.11.2. et}{insertion}{ega25_25_181-v0019} achève la démonstration de 16.9.

NB On aurait pû donner 16.11.2. dès le début en corollaire à
16.2. ou 16.3. \hand{(} qui auraient dû s'interchanger); ainsi la partie du n\textsuperscript{o}
indépendante du ``discriminant corrigé'' serait bloquée au début du n\textsuperscript{o}.
D'ailleurs tu dois mieux savoir que moi dans quelle mesure la notion
de discriminant corrigé et 16.9. étai\unsure{e|t}nt connus, pour en créditer qui de
droit. Peut-être y a-t\unsure{-}il déjà une autre terminologie reçue ? Je remarque
d'ailleurs que c'est le discriminant \uline{\unsure{n|c}o\unsure{d|r}rigé} qui mérité vraiment le nom
de discriminant, c'est lui qui se généralise en le discriminant d'une
forme quelconque (cf N\textsuperscript{o} suivant). Le discriminant d'une forme quadratique
(dans la terminologie adoptée ici, et qui, si je ne me trompe, est la

\mspage{\hand{181}}
